\documentclass[11pt, oneside]{article}   	% use "amsart" instead of "article" for AMSLaTeX format
\usepackage{geometry}                		% See geometry.pdf to learn the layout options. There are lots.
\geometry{a4paper}                   		% ... or a4paper or a5paper or ... 
%\geometry{landscape}                		% Activate for rotated page geometry
%\usepackage[parfill]{parskip}    		% Activate to begin paragraphs with an empty line rather than an indent
\usepackage{graphicx}				% Use pdf, png, jpg, or eps§ with pdflatex; use eps in DVI mode
								% TeX will automatically convert eps --> pdf in pdflatex		
\usepackage{amssymb}
\usepackage{amsmath}
\usepackage{authblk}
\usepackage[hyphens]{url}

\usepackage{url}
\usepackage{hyperref} 

\usepackage{titlesec}

\usepackage{caption}


\title{Remote Energy Monitoring}
\author{BEng(Hons) MSc PhD(Cantab) GMCOP CEng MCIBSE}

\date{\today}							% Activate to display a given date or no date

\begin{document}
\maketitle
\begin{center}
\href{mailto:aidan.parkinson@gmail.com}{aidan.parkinson@gmail.com}
\end{center}

Many organisations intend to upgrade to a remote energy monitoring system, to understand energy use and improve service performance across their globally distributed estates.
Having a better understanding of energy use should allow strategies and programmes to be developed holistically at a local level.
The following requirements provide guidance on how to inter-network with such an application.

\section{Utility Metering}

The tables below compare the various types of utility meters available.

\begin{table}[H]
	\caption{Internet-of-Things (IoT) meter}
	\begin{center}
	\begin{tabular}{| l | l |}
	\hline
	Advantages&Disadvantages\\
	\hline
	No requirement for inter-networking gateway&Expensive
	Meter counts energy consumption&Device testing requirement to validate performance
	Potentially more descriptive information payload.	
	\hline
	\end{tabular}
	\end{center}
	\label{IoT Meter Option}
\end{table}

\begin{table}[H]
	\caption{Open building systems communications protocol meter}
	\begin{center}
	\begin{tabular}{| l | l |}
	\hline
	Advantages&Disadvantages\\
	\hline
	Meter counts energy consumption&Insufficient security features for exposure to the Wide Area Network (WAN).
	Potentially more descriptive information payload.	
	\hline
	\end{tabular}
	\end{center}
	\label{Open Protocol Meter Option}
\end{table}

\begin{table}[H]
	\caption{Pulse meter}
	\begin{center}
	\begin{tabular}{| l | l |}
	\hline
	Advantages&Disadvantages\\
	\hline
	Relatively cheap&Requires a pulse counting device
	&Minimal information payload
	\hline
	\end{tabular}
	\end{center}
	\label{Pulse Meter Option}
\end{table}

\section{Inter-networking}

It is important that any inter-networking gateways be hosted on a dedicated Virtual Private Network, with layer 3 separation from other building management system networks.
Otherwise, exposure of such inter-networking gateways to the Wide Area Network would pose an unacceptable cybersecurity risk.

\begin{table}[H]
	\caption{Inter-Networking Gateway}
	\begin{center}
	\begin{tabular}{| l | l |}
	\hline
	Inter-Networking Standard&Justification\\
	\hline
	802.1x&Certificate based network access managed by a Remote Authentication Dial-In User Service (RADIUS)
	IP version 6&Unique, lower cost and permanent public Internet Protocol (IP) addressing, with native support for IPSec encryption of entire data packets in transit.
	TLS version 1.3&Most recent release of Transport Layer Security (TLS) protocol, for encryption of data packets transport layer in transit.
	LDAP&Lightweight Directory Access Protocol (LDAP) is used for accessing and maintaining distributed directory information services over an Internet Protocol (IP) network. A common use of LDAP is to provide a central place to store usernames and passwords.
	\hline
	\end{tabular}
	\end{center}
	\label{Pulse Meter Option}
\end{table}

hoshi.property prefers telemetry over Secure Message Queue Telemetry Transport (MQTTS) for the following reasons:

\item[ a)] All inter-networking communications managed from a single endpoint;
\item[ b)] Lightweight telemetry;
\item[ c)] No packet loss, even through interruptions to network connectivity.

A MQTTS message broker is available at the URL: mqtts.hoshi.property and may subscribe to secure telemetry, given recognised 802.1X and TLS certificates.

Each inter-networking gateway shall be allocated a Globally Unique Identifier (GUID), serving as the message brokers topic.
Customers may link available Guid's to service addresses through the hoshi.property user interface.

Each topic requires the following response payload schema:

"$schema": "https://json-schema.org/draft/2020-12/schema",
    "$id": "https://example.com/product.schema.json",
    "title": "Utilisation",
    "description": "Utilisation telemetry",
    "type": "object",
    "properties": {
        "required": ["Guid", "utilisation", "time"],
        "Guid": {
            "description": "Globally Unique Identifier inter-networking gateway",
            "type": "string"
        },
		"time": {
			"desription" : "Universal time telemetry was published to mqtts.hoshi.property"
			"value" : setUTCDate()
		}
		"Utilisation": {
			{
            "description": "Electric branch meter",
			"topic": "EBM",
			"value": number
		},
		{
            "description": "Electricity meter",
			"topic": "EM",
			"value": number
		},
		{
            "description": "Virtual electricity meter",
			"topic": "EMV",
			"value": number
		},
		{
            "description": "Flow meter",
			"topic": "FM",
			"value": number
		},
		{
            "description": "Gas meter",
			"topic": "GM",
			"value": number
		},
		{
            "description": "Virtual gas meter",
			"topic": "GMV",
			"value": number
		},
		{
            "description": "Heat meter",
			"topic": "HM",
			"value": number
		},
		{
            "description": "virtal heat meter",
			"topic": "HMV",
			"value": number
		},
		{
            "description": "water meter",
			"topic": "WM",
			"value": number
		},
		{
            "description": "virtual water meter",
			"topic": "WMV",
			"value": number
		}
	}}

\end{document}